\begin{claim}{C.47}{Proved}{Three independent boundary quotients force a Lie action}
Let $\Phi_i:O\to B_i$, $1\leq i\leq3$, be three actual boundary
quotients of the same open orbit as in C.42. Use their common source
universal covering and real target coverings, with
\[
 p:X\longrightarrow O,\qquad
 q_iF_i=\Phi_i p,\qquad F_i:X\longrightarrow\R .
                                                        \tag{C.47.1}
\]
Suppose the three lifted fiber partitions are pairwise distinct;
equivalently, none of the pairs satisfies the first alternative
of C.45. The full group
\[
 \mathcal W=
 \{a\in\operatorname{Homeo}(X):
       F_i a=b_iF_i\text{ for some }
       b_i\in\operatorname{Homeo}^+(\R),\ 1\leq i\leq3\}
                                                        \tag{C.47.2}
\]
is a closed, locally compact subgroup in the actual compact-open
topology. It is a Lie group of dimension two or three and acts
transitively on $X$. If any pair $(F_i,F_j)$ has image different
from $\R^2$, this action is simply transitive and has dimension two.
In particular, the closure in $\operatorname{Homeo}(X)$ of the
compatible lifted loop group is locally compact.
All three covering equations in (C.47.1) and the kernel of the
global-to-restricted action are retained. The conclusion concerns
the action on the source covering plane; it does not assert an
extension of an intrinsic homotopy across the frontier of $O$.
\end{claim}

\begin{proof}
\emph{The actual group topology.}
Include all three target lifts in the single closed Polish group
of tuples $(l,a,b_1,b_2,b_3)$ satisfying
\[
 pa=(l|_O)p,\qquad q_i b_i=(l|_{B_i})q_i,\qquad
 F_i a=b_iF_i \quad(1\leq i\leq3).
                                                        \tag{C.47.3}
\]
For each source lift, the three compatible target lifts are
obtained exactly as in C.42--C.45. This group acts transitively
on $X$, so its image in $\mathcal W$ does also.
Each $b_i$ in (C.47.2) is unique because $F_i$ is onto.
C.43 supplies a continuous section
$\sigma_i:\R\to X$ of each $F_i$. Consequently
\[
              b_i=F_i a\sigma_i,\qquad
              b_i^{-1}=F_i a^{-1}\sigma_i.               \tag{C.47.4}
\]
These formulas are continuous for uniform convergence on compact
sets, for maps and inverses.

They also prove closedness. If $a_n\in\mathcal W$ converges to a
homeomorphism $a$ together with its inverse, (C.47.4) gives locally
uniform limits $b_i$ and $d_i$ of the target maps and their
inverses. The identities $b_{i,n}d_{i,n}=d_{i,n}b_{i,n}=\id$
pass to the limit, since the images of each compact set remain
in a common compact set. Thus $d_i=b_i^{-1}$; each $b_i$ is an
increasing homeomorphism, and $F_i a=b_iF_i$. Hence $\mathcal W$
is a closed Polish subgroup of $\operatorname{Homeo}(X)$.

C.45 identifies $X$ by any two of the $F_i$ with an open domain
having interval coordinate sections. In coordinates $(x,y)=(F_1,F_2)$,
write $F_3=w(x,y)$. Every restriction of $w$ to a horizontal or
vertical section is strictly monotone. The signs are constant
over the connected domain: for two fixed nearby values in one
coordinate, the continuous difference of their $w$-values cannot
change sign. The coordinate representation of $\mathcal W$ is
therefore exactly the triples $(g_1,g_2,k)$ satisfying
\[
 (g_1,g_2)\Omega=\Omega,\qquad
 w(g_1x,g_2y)=k(w(x,y)).                                 \tag{C.47.5}
\]
The map from $\mathcal W$ to these triples and its inverse are
continuous, by (C.47.4) and the pair-coordinate homeomorphism.
All topological group assertions below use this topology.

\emph{A strictly interior mean makes an ordered-pair stabilizer trivial.}
We first prove an elementary fact used in the non-plane cases.
Let
\[
 D=\{(a,b):a<b\}
 \quad\hbox{or}\quad
 D=\{(a,b):a<b<a+1\}.
                                                        \tag{C.47.6}
\]
Suppose a group of increasing line homeomorphisms preserves $D$
diagonally and preserves a continuous function $M:D\to\R$ with
\[
 a<M(a,b)<b,
 \qquad M\text{ injective in each argument separately}.   \tag{C.47.7}
\]
A group element fixing one pair $(a,b)\in D$ is the identity.
Indeed, it fixes all the points obtained by repeatedly inserting
$M$ between consecutive fixed endpoints. These subdivisions
have mesh tending to zero. To verify this, fix $\varepsilon>0$.
On the compact set of pairs
\[
 a\leq u<v\leq b,\qquad v-u\geq\varepsilon,
 \]
both $M(u,v)-u$ and $v-M(u,v)$ have a positive minimum.
Every such pair is still in $D$, including in the band case
because $v\leq b<a+1\leq u+1$.
Every subdivision of an interval of length at least
$\varepsilon$ therefore shortens each child by a fixed positive
amount. No sufficiently deep subdivision can have that length.
Continuity now makes the homeomorphism fix $[a,b]$ pointwise.

Let $I$ be the maximal interval of fixed points containing
$[a,b]$. A finite right endpoint $c$ belongs to the fixed set.
Choose $x<c$ in the interior of $I$ with $(x,c)\in D$.
For $y>c$ sufficiently close to $c$, the pair $(x,y)$ remains in
$D$ and $M(x,y)$ remains in the interior of $I$. Equivariance
gives
\[
 M(x,gy)=M(gx,gy)=gM(x,y)=M(x,y).
 \]
Injectivity in the second argument forces $gy=y$, extending
$I$ past $c$. The analogous argument, with a fixed $y>c$ and
$x<c$ close to a finite left endpoint, uses injectivity in the
first argument. Both finite endpoints are impossible.
Thus $I=\R$ and $g=\id$.

\emph{A non-plane pair and two full remaining pairs.}
Suppose some pair image is not $\R^2$. Relabel and apply the
explicit coordinate changes of C.45 so that
$\Omega=D$ is one of (C.47.6).
The first two target actions then coincide:
$g_1=g_2=g$; in the band case $g(x+1)=g(x)+1$.
First suppose the images of $(F_1,F_3)$ and $(F_2,F_3)$ are both
the full plane. Each $w$-level is then the graph
\[
                   y=r_z(x),\qquad z\in\R,               \tag{C.47.8}
\]
where $(x,z)\mapsto r_z(x)$ is continuous and every $r_z$ is
a homeomorphism of $\R$ onto $\R$.
It is increasing, since $r_z(x)>x$ for all $x$, which a
decreasing real-line homeomorphism cannot satisfy.
Reverse the $z$-coordinate when necessary. Then
$z\mapsto r_z(x)$ is strictly increasing, and its range is
$(x,+\infty)$ in the half-plane case and $(x,x+1)$ in the
band case. The sign is the same for every $x$ by connectedness.
The action equation is exactly
\[
                      g r_z g^{-1}=r_{kz}.               \tag{C.47.9}
\]

For $(a,b)\in D$, there is a unique $s=s(a,b)$ with
\[
                   r_s^2(a)=b.
 \quad\hbox{Define}\quad M(a,b)=r_s(a).                  \tag{C.47.10}
\]
Here are the endpoint and continuity checks. The left side is
continuous and strictly increasing in $s$.
As $s\to-\infty$, $r_s(x)\to x$, and hence $r_s^2(a)\to a$.
This last limit follows by bounding $r_s(a)$ between $a$ and
$a+\delta$ and then using monotonicity of $r_s$ at these two
fixed points. As $s\to+\infty$ the square tends to $+\infty$
in the half-plane case. In the band case $r_s(x)\to x+1$,
and the same fixed-point bracketing gives
$r_s^2(a)\to a+2$. The required value $b$ lies strictly between
the endpoint limits. Strict monotonicity gives uniqueness,
and two nearby $s$-values bracketing a solution continue to
bracket it when $(a,b)$ varies slightly; this proves continuity.

The value in (C.47.10) lies strictly between $a$ and $b$.
It is strictly increasing in $b$ for fixed $a$.
For fixed $b$ and $a_1<a_2$, the corresponding parameters satisfy
$s_1>s_2$, because $r_s^2$ increases in its argument.
If $r_{s_1}(a_1)\geq r_{s_2}(a_2)$, applying these increasing
maps once more and using $s_1>s_2$ gives
\[
 b=r_{s_1}(r_{s_1}(a_1))
   >r_{s_2}(r_{s_2}(a_2))=b,
 \]
a contradiction. Thus $M$ is strictly increasing in its first
argument as well.
Equation (C.47.9) and uniqueness in (C.47.10) prove
$M(ga,gb)=gM(a,b)$.
The preceding mean argument therefore makes every source
point stabilizer trivial: once $g=\id$, surjectivity of $w$
in (C.47.5) also gives $k=\id$.

\emph{A further non-plane pair identifies all three line actions.}
It remains, with $\Omega=D$, to consider the case where at least
one of the other pair images is not the full plane. One of its
finite endpoint homeomorphisms from C.45 gives an exact equation
\[
                         f g=k f.                       \tag{C.47.11}
\]
Replace the third coordinate by $f^{-1}F_3$.
Now all three target actions are the same $g$; this remains true
when $f$ is decreasing. Denote the new third coordinate again
by $w(x,y)$. An equality $w=x$ or $w=y$ at one point would
hold everywhere by transitivity and contradict independence.
The relative order of the three values is therefore constant.

In the half-plane, the projected group of $g$'s acts transitively
on all ordered pairs. Its centralizer in
$\operatorname{Homeo}(\R)$ is trivial: the stabilizer of a point
$x$ acts transitively on each of $(-\infty,x)$ and $(x,+\infty)$,
so its only common fixed point is $x$.
A commuting map must send $x$ to that same fixed point.
It follows that every finite endpoint of either remaining
pair domain, now an equivariant self-homeomorphism of the
same line, is the identity.

When $x<w(x,y)<y$, take $M=w$.
When $w(x,y)>y$, the $(x,w)$ image is exactly
$\{(x,z):x<z\}$: its finite lower endpoint must be the identity,
and a finite upper endpoint would be the same map and leave
an empty interval. For each $a<b$, define $M(a,b)=c$ by
\[
                         w(a,c)=b.                      \tag{C.47.12}
\]
The restriction in $c$ is increasing and maps $(a,+\infty)$
onto $(a,+\infty)$; its increase also follows from
$w(a,c)>c$ and the limit as $c\to+\infty$.
The unique solution satisfies $a<c<b$.
When $w(x,y)<x$, the $(y,w)$ image is exactly
$\{(y,z):z<y\}$; define instead
\[
                         w(c,b)=a.                      \tag{C.47.13}
\]
The restriction in $c$ is increasing and maps
$(-\infty,b)$ onto $(-\infty,b)$, so again $a<c<b$.
Both inverse constructions are continuous by the pair-coordinate
homeomorphisms. They are injective in each argument, by the
strict injectivity of the corresponding restrictions of $w$.
Equation (C.47.5), with all three actions equal, proves their
equivariance.

In the band, $g$ commutes with unit translation.
The equalities $w=x+n$ and $w=y+n$, for any integer $n$, are
excluded by the same transitivity and independence argument.
There is consequently one integer $n$, constant on the connected
domain, such that
\[
                         x<w(x,y)-n<x+1.
 \]
Either $w(x,y)-n<y$ everywhere, or the reverse inequality holds
everywhere. In the first case take $M(x,y)=w(x,y)-n$.
In the second case use the exact domain homeomorphism
\[
          (u,v)=(y,x+1),\qquad
          M(u,v)=w(v-1,u)-n.                            \tag{C.47.14}
\]
It sends the band onto itself, and $u<M(u,v)<v$.
Both constructions are continuous and injective in each argument.
They are equivariant because every $g$ commutes with the
integer shifts. We have obtained a mean satisfying
(C.47.7) in all cases.

The mean argument proves simple transitivity of $\mathcal W$
whenever a pair image is not the full plane.
The Polish open orbit-map theorem makes evaluation
$\mathcal W\to X$ a homeomorphism in this case.
Thus $\mathcal W$ is locally compact and is a topological
two-manifold. Its faithful product action on $X\times\R$
and the Hilbert--Smith theorem in dimension three
\cite[Theorem~1.5 and the introductory reduction]{C_Pardon2013}
make it a two-dimensional Lie group.

\emph{Three full pair images give a global ordered loop.}
Suppose now that all three pair images are $\R^2$.
The function $w:\R^2\to\R$ is then continuous, and each partial
map is a homeomorphism onto $\R$. Reverse input coordinates
as necessary to make it increasing in both arguments.
At a chosen source point $(a,b)$ put $c=w(a,b)$ and use the
explicit coordinates
\[
 u=w(x,b)-c,\qquad v=w(a,y)-c,\qquad z=w(x,y)-c.
                                                        \tag{C.47.15}
\]
They give a continuous multiplication $m:\R^2\to\R$ with
\[
 m(u,0)=u,\qquad m(0,v)=v,                              \tag{C.47.16}
\]
strictly increasing in both variables, each partial map onto
$\R$. Its two division maps are continuous: solving a strictly
increasing continuous equation is continuous by fixed
bracketing on either side of the solution.

Let $S$ be the stabilizer of $(0,0)$ in $\mathcal W$.
Setting either argument equal to zero in its action equation
shows that the three target maps coincide. Thus, in its actual
subspace topology, $S$ is exactly
\[
 \operatorname{Aut}^+(m)=
 \{f\in\operatorname{Homeo}^+(\R):
       f(0)=0,\quad f(m(x,y))=m(fx,fy)\}.                 \tag{C.47.17}
\]
This is a closed subgroup. We next prove its local compactness
by uniform estimates for both maps and inverses.

Fix $a>0$ and $0<A<B$. Set
\[
 \mathcal F_{A,B}=\{f\in\operatorname{Aut}^+(m):
                            A\leq f(a)\leq B\}.
                                                        \tag{C.47.18}
\]
Define $D(x)=m(x,x)$ and $T_c(x)=m(c,x)$.
The map $D$ is an increasing homeomorphism,
$D(x)>x$ for $x>0$, and $D(x)<x$ for $x<0$.
The numbers
\[
 v_n=D^{-n}(a)>0,\qquad
 u_n\text{ defined by }m(u_n,v_n)=0                      \tag{C.47.19}
\]
therefore converge to zero, with $u_n<0$.
Every automorphism respects these formulas, and hence
\[
 D^{-n}(A)\leq f(v_n)\leq D^{-n}(B).
                                                        \tag{C.47.20}
\]
The two bounds tend to zero. If $\ell(t)$ is the unique
solution of $m(\ell(t),t)=0$, then
$f(u_n)=\ell(f(v_n))$; continuity of $\ell$ makes this tend
to zero uniformly over $\mathcal F_{A,B}$ as well.

For $c>0$, $T_c$ is increasing, onto, and satisfies $T_c(x)>x$.
Its integer iterates of zero tend to both ends of $\R$.
For every integer $j$,
\[
                       f(T_a^j(0))=T_{f(a)}^j(0).       \tag{C.47.21}
\]
To bound $f$ on a given compact interval, enclose it between
two such iterates with $j=-N,N$.
The right sides of (C.47.21) range over a compact set as
$f(a)$ ranges over $[A,B]$; the negative iterates depend
continuously on $c$ by the division-map continuity.
Monotonicity then bounds all $f$ uniformly on that interval.

For equicontinuity, let $I$ be a compact interval.
For fixed $n$, both gaps in
\[
                 m(u_n,x)<x<m(v_n,x)\qquad(x\in I)
 \]
have positive minima on $I$. Thus a single $\delta_n>0$
ensures
\[
 |y-x|<\delta_n,\ x\in I
 \quad\Longrightarrow\quad
                   m(u_n,x)<y<m(v_n,x).
 \]
Applying $f$ bounds $f(y)$ between
$m(f(u_n),f(x))$ and $m(f(v_n),f(x))$.
The already proved compact bound on $f(I)$, the uniform limits
$f(u_n),f(v_n)\to0$, and continuity of $m$ show that both
bounds approach $f(x)$ uniformly in $x$ and $f$.
This is equicontinuity on every compact interval.

The inverse family has the same estimates with new fixed
positive bounds. Choose $n$ so large that $D^{-n}(B)<a$.
Then $f(v_n)<a$ and $f^{-1}(a)>v_n>0$.
Choose $N$ with $T_A^N(0)>a$. Since $T_c^N(0)$ increases
with $c>0$, (C.47.21) gives
$f(T_a^N(0))\geq T_A^N(0)>a$, and hence
\[
                     v_n<f^{-1}(a)<T_a^N(0).            \tag{C.47.22}
\]
Applying the same compact bounds and equicontinuity argument
to these positive bounds controls all inverses uniformly.

Arzel\`a--Ascoli on a countable compact exhaustion now gives,
from any sequence in $\mathcal F_{A,B}$, locally uniform
limits $f$ and $h$ of a common subsequence of maps and inverses.
Their compact image bounds permit passage to the composition
identities, so $fh=hf=\id$.
The multiplication equation and the bounds in (C.47.18)
pass to the limit. Thus $\mathcal F_{A,B}$ is compact in the
homeomorphism-group topology.
Taking $A<a<B$ gives a compact neighborhood of the identity
in $S$, and proves its local compactness.

Evaluation $f\mapsto f(a)$ is also injective.
An automorphism fixing $a$ fixes every $v_n$ and therefore every
$T_{v_n}^j(0)$, for $j\in\Z$.
For fixed $x$, choose the first such iterate at or to the right
of $x$. It lies in $[x,T_{v_n}(x)]$, whose length tends to
zero since $v_n\to0$. These fixed points tend to $x$, so
the automorphism fixes every $x$.
The compactness of (C.47.18) proves that this continuous
evaluation is a proper embedding
\[
                              S\lhook\joinrel\longrightarrow(0,\infty).
                                                        \tag{C.47.23}
\]

\emph{From the stabilizer to the full group, without changing topology.}
The locally compact group $S$ acts faithfully and continuously
on $\R^3$ by its action on the first coordinate.
Pardon's theorem makes $S$ a Lie group.
Its embedding (C.47.23) gives $\dim S\leq1$.
The full group $\mathcal W$ is still the closed Polish group
already defined; it has not yet been assumed locally compact.

Consider the free right action of $S$ on $\mathcal W$.
It is a Cartan action in its actual topology.
Indeed, choose an identity neighborhood $U$ in $\mathcal W$
such that $U\cap S$ has compact closure in $S$, using the
proved local compactness and the subspace topology of $S$.
Choose an identity neighborhood $V$ with $V^{-1}V\subset U$.
Then
\[
       \{s\in S:Vs\cap V\ne\varnothing\}
                   \subseteq V^{-1}V\cap S
 \]
has compact closure in $S$, and translates provide the
required neighborhoods at every point.
The space $\mathcal W$ is metrizable and completely regular.
The free-action form of Palais' slice theorem therefore gives
local sections and local trivializations for
\[
                  \mathcal W\longrightarrow\mathcal W/S
 \quad\cite[Theorem~2.3.3]{C_Palais1961}.
 \]
For clarity, the slice map has the form
$\sigma(u)s\mapsto(u,s)$; its inverse is multiplication.
If $\theta$ is the equivariant slice function to $S$, the
inverse coordinates of $w$ are
$(w\theta(w)^{-1},\theta(w))$, as in the explicit free-action
statement \cite[Theorem~3.7]{C_KramerMuratGarcia2026}.
The open orbit-map theorem identifies
$\mathcal W/S$ homeomorphically with $X=\R^2$.
These actual local products of the locally compact spaces
$S$ and an open subset of $X$ prove that $\mathcal W$ is
locally compact. Its faithful action on $X\times\R$ then
makes it a Lie group by Pardon again.
The same local products give
$\dim\mathcal W=2+\dim S$, which is two or three.

Finally, every compatible lifted loop value belongs to
$\mathcal W$, and its closure in $\operatorname{Homeo}(X)$
is a closed subgroup of this locally compact group.
It is therefore locally compact with precisely the induced
homeomorphism-group topology.
None of the arguments replaces the $l$-coordinate of
(C.47.3) by its action on $X$: the kernel of that restriction
can still act elsewhere.
The coordinate changes act on the specified real target lifts
and leave the original equations (C.47.1) available throughout.
This proves the asserted intrinsic reduction and leaves the
frontier extension step explicit.
\end{proof}
